\section{Experiment}

    \subsection{Versuchsaufbau}
        Die Messung der Hysteresekurven, also die magnetische Flussdichte $B$ in Abhängigkeit des äußeren Magnetfeldes $H$,
        erfolgt mit einer Hysteresemessbrücke, welche in Abbildung \ref{fig:Versuchsaufbau} dargestellt ist.
        \begin{figure}[H]
            \centering
            \includegraphics[width=0.4\textwidth]{Versuchsaufbau.png}
            \caption{Versuchsaufbau zur Messung der Hysteresekurven \cite{Versuchsaufbau}.}
            \label{fig:Versuchsaufbau}
        \end{figure}
        Am X-Eingang wird die Spannung $U_f$ über ein Widerstand R abgegriffen. Diese Spannung ist proportional zur
        Feldstärke $H$, welche sich mit dem Ampere'schen Gesetz berechnen lässt:
        \begin{equation}
            \int H ds = N_1 I \implies H = \frac{N_1 I}{\pi d R}
            \label{eq:Feldstaerke}
        \end{equation}
